\section{中心胞元约化与连续 DtN 算子}
\label{sec:dtn}

以下固定第~\ref{sec:modelproblem}~节中的 $\beta$ 以及选定谱隙
$G_\beta=(g_-,g_+)$。沿 $y$ 方向施加 $\beta$-准周期条件以后，条带
$B=\RR\times(-d/2,d/2)$ 在 $x$ 方向仍然无界。记
$I_d=(-d/2,d/2)$。由式~\eqref{eq:defect-region} 的全平面宏观分区定义
\begin{align*}
  B_\varsigma&=\mathcal{W}_\varsigma\cap B,
  \qquad \varsigma\in\{-,0,+\},\\
  B_-&=(-\infty,-a)\times I_d,
  &B_0&=(-a,a)\times I_d,
  &B_+&=(a,\infty)\times I_d.
\end{align*}
中心缺陷胞元及紧邻它的左右周期胞元分别定义为
\begin{align}
  \mathcal{C}_0&=B_0=(-a,a)\times I_d,
  &\mathcal{C}_-&=(-a-L,-a)\times I_d,
  &\mathcal{C}_+&=(a,a+L)\times I_d.
  \label{eq:center-and-lead-cells}
\end{align}
这里 $\mathcal{C}_0$ 是宽度为 $2a$ 的功能性缺陷胞元，其宽度不要求等于 $L$；
$\mathcal{C}_-$ 和 $\mathcal{C}_+$ 才是 $x$ 方向周期为 $L$ 的一个基本背景胞元，而不是
半无限区域。记 $e_x=(1,0)$ 为 $x$ 方向的单位向量。左右半条带可以由这些周期胞元表示为
\begin{align}
  B_-
  &=\operatorname{int}\!\left(
    \bigcup_{j=0}^{\infty}
    (\overline{\mathcal{C}_-}-jL e_x)\right),
  &
  B_+
  &=\operatorname{int}\!\left(
    \bigcup_{j=0}^{\infty}
    (\overline{\mathcal{C}_+}+jL e_x)\right).
  \label{eq:half-guides}
\end{align}
对 $\varsigma\in\{-,0,+\}$，以
$\Omega_\varsigma\subset\mathcal{C}_\varsigma$ 表示相应胞元中的单个介质柱截面区域。
其中 $\Omega_0$ 是中心缺陷胞元中的介质柱区域，并允许
$\Omega_0=\varnothing$；$\Omega_-$ 和 $\Omega_+$ 分别是左右基本周期胞元中的介质柱区域。
中心介质柱若存在则完全包含在 $\mathcal{C}_0=B_0$ 中；$B_-$ 和 $B_+$ 内的介质柱区域分别由
$\Omega_-$ 的平移 $\Omega_--jLe_x$ 和 $\Omega_+$ 的平移 $\Omega_++jLe_x$，
$j\geq0$，生成。因此 $B_\pm$ 中的系数均为周期背景系数 $\rho_{\mathrm p}$，
$\mathcal{C}_0$ 中则保留第~\ref{sec:modelproblem}~节定义的缺陷系数 $\rho$ 及
$\Upsilon$ 所表示的全部尖锐材料界面。人工界面和法向约定为
\begin{align}
  \Gamma_\pm
  &=\partial\mathcal{C}_0\cap\partial B_\pm,
  &\Gamma_-&=\{-a\}\times I_d,
  &\Gamma_+&=\{a\}\times I_d,
  \notag\\
  \nu_-&=-e_x,
  &\nu_+&=e_x.
  \label{eq:artificial-interfaces}
\end{align}
两个法向都从 $\mathcal{C}_0$ 指向相应半条带 $B_\pm$。特别地，它们是中心区域的外法向，
却与 $B_\pm$ 在 $\Gamma_\pm$ 上的外法向相反；下文的 $\partial_{\nu_\pm}$ 始终采用
式~\eqref{eq:artificial-interfaces} 的定向。

对 $U\in\{\mathcal{C}_0,B_-,B_+\}$，只引入后文所需的准周期能量空间
\begin{align}
  H^1_\beta(U)
  :=\left\{z\in H^1(U):
  z|_{y=d/2}=\ee^{\ii\beta d}z|_{y=-d/2}
  \text{，按迹意义}\right\}.
  \label{eq:local-quasiperiodic-space}
\end{align}
令 $\gamma_\pm$ 表示 $\Gamma_\pm$ 上的 Dirichlet 迹，并以迹像及其商范数定义
\begin{align}
  H^{1/2}_\beta(\Gamma_\pm)
  &:=\gamma_\pm H^1_\beta(\mathcal{C}_0)
    =\gamma_\pm H^1_\beta(B_\pm),
  &
  H^{-1/2}_\beta(\Gamma_\pm)
  &:=H^{1/2}_\beta(\Gamma_\pm)^*.
  \label{eq:interface-trace-spaces}
\end{align}
这样，$y$ 方向的 $\beta$-准周期兼容性已经编码在迹空间中，无须借助临界
$H^{1/2}$ 函数的端点点值。记
$\langle\cdot,\cdot\rangle_{\Gamma_\pm}$ 为
$H^{-1/2}_\beta(\Gamma_\pm)$ 与 $H^{1/2}_\beta(\Gamma_\pm)$ 之间的对偶配对。

在每条材料界面上任意固定一个两侧共用的单位法向 $\nu$。现在固定 $k>0$ 且
$k^2\in G_\beta$。给定任一方向和界面数据
$\varphi\in H^{1/2}_\beta(\Gamma_\pm)$，相应的连续半波导 Dirichlet 问题是求
$u_\pm=u_\pm(\varphi;k)\in H^1_\beta(B_\pm)$，使
\begin{equation}
  \left\{
  \begin{aligned}
    -\Delta u_\pm&=k^2\rho_{\mathrm p}u_\pm
    &&\text{在 }B_\pm\setminus\Upsilon\text{ 的各均匀材料子区域内},\\
    u_\pm^+&=u_\pm^- ,
    &\partial_\nu u_\pm^+&=\partial_\nu u_\pm^-
    &&\text{在 }\Upsilon\cap B_\pm,\\
    u_\pm|_{y=d/2}
    &=\ee^{\ii\beta d}u_\pm|_{y=-d/2},
    &\partial_yu_\pm|_{y=d/2}
    &=\ee^{\ii\beta d}\partial_yu_\pm|_{y=-d/2},\\
    u_\pm|_{\Gamma_\pm}&=\varphi.&&
  \end{aligned}
  \right.
  \label{eq:half-guide-dirichlet-problem}
\end{equation}
第二行正是当前分片常数模型的传输条件。解还必须属于谱隙中向相应无穷远衰减的分支。
具体而言，存在与 $k$ 有关的
$c_{\mathrm{dec},\pm},\alpha_\pm>0$，使对 $j\geq0$ 有
\begin{align}
  \|u_-\|_{H^1(\mathcal{C}_--jLe_x)}
  &\leq c_{\mathrm{dec},-}\ee^{-\alpha_-j}
    \|\varphi\|_{H^{1/2}_\beta(\Gamma_-)},
  &
  \|u_+\|_{H^1(\mathcal{C}_++jLe_x)}
  &\leq c_{\mathrm{dec},+}\ee^{-\alpha_+j}
    \|\varphi\|_{H^{1/2}_\beta(\Gamma_+)}.
  \label{eq:half-guide-decay}
\end{align}
这个条件比含糊的辐射条件更强，并与导模在 $B$ 中的有限能量以及所选投影谱隙一致。

谱隙条件本身并不保证问题~\eqref{eq:half-guide-dirichlet-problem} 唯一可解。记
$\Sigma_{D,\pm}(\beta)\subset G_\beta$ 为在 $\Gamma_\pm$ 施加齐次 Dirichlet 条件后，
相应半条带 $B_\pm$ 上仍存在非零衰减解的例外谱。Fliss 的连续模型在本质谱外除一个可数例外集外给出
半波导问题的适定性\cite[定理~4.1]{Fliss2013}；在当前分片常数模型中，这一例外集必须针对
包含传输条件的半波导算子定义，不能从连续折射率模型直接照搬。下文只考虑
\begin{align}
  k^2\in G_\beta\setminus
  \bigl(\Sigma_{D,-}(\beta)\cup\Sigma_{D,+}(\beta)\bigr),
  \label{eq:dtn-admissible-parameters}
\end{align}
并要求问题~\eqref{eq:half-guide-dirichlet-problem} 对每个 $\varphi$ 唯一可解。

在参数条件~\eqref{eq:dtn-admissible-parameters} 下，定义连续 DtN 算子
\begin{align}
  \Lambda_\pm(k):H^{1/2}_\beta(\Gamma_\pm)
  &\longrightarrow H^{-1/2}_\beta(\Gamma_\pm),
  &
  \Lambda_\pm(k)\varphi
  &:=\left.\partial_{\nu_\pm}u_\pm(\varphi;k)
    \right|_{\Gamma_\pm}.
  \label{eq:continuous-dtn}
\end{align}
其中法向导数按弱共法向迹理解。由于 $\nu_\pm$ 从中心指向 $B_\pm$，半条带自身外法向下的
导数等于 $-\Lambda_\pm(k)\varphi$。在 Dirichlet 例外参数处，单值算子
\eqref{eq:continuous-dtn} 可能失效；更稳健的连续对象是所有衰减半波导解的 Dirichlet--Neumann
Cauchy 数据所组成的线性关系。只有该关系到 Dirichlet 分量的投影可逆时，它才是
$\Lambda_\pm(k)$ 的图。本节始终处在式~\eqref{eq:dtn-admissible-parameters} 的单值情形。

利用这两个算子，可以把无界问题限制到 $\mathcal{C}_0$。中心场
$u_0\in H^1_\beta(\mathcal{C}_0)\setminus\{0\}$ 应满足
\begin{equation}
  \left\{
  \begin{aligned}
    -\Delta u_0&=k^2\rho u_0
    &&\text{在 }\mathcal{C}_0\setminus\Upsilon\text{ 的各均匀材料子区域内},\\
    u_0^+&=u_0^- ,
    &\partial_\nu u_0^+&=\partial_\nu u_0^-
    &&\text{在 }\Upsilon\cap\mathcal{C}_0,\\
    u_0|_{y=d/2}
    &=\ee^{\ii\beta d}u_0|_{y=-d/2},
    &\partial_yu_0|_{y=d/2}
    &=\ee^{\ii\beta d}\partial_yu_0|_{y=-d/2},\\
    \partial_{\nu_\pm}u_0
    &=\Lambda_\pm(k)\gamma_\pm u_0
    &&\text{在 }\Gamma_\pm.
  \end{aligned}
  \right.
  \label{eq:center-dtn-problem}
\end{equation}
最后一行是精确透明边界条件。对任意 $v\in H^1_\beta(\mathcal{C}_0)$ 分部积分时，顶、底
边界项因准周期条件相消，而 $\Gamma_\pm$ 的中心外法向恰为 $\nu_\pm$。代入
式~\eqref{eq:center-dtn-problem} 的最后一行，得到弱形式
\begin{align}
  \int_{\mathcal{C}_0}\nabla u_0\mathbin{\cdot}\nabla\overline v
  -k^2\int_{\mathcal{C}_0}\rho u_0\overline v
  -\sum_{\varsigma\in\{-,+\}}
   \left\langle
   \Lambda_\varsigma(k)\gamma_\varsigma u_0,
   \gamma_\varsigma v
   \right\rangle_{\Gamma_\varsigma}
  &=0.
  \label{eq:center-dtn-weak-form}
\end{align}
两个 DtN 项的负号都来自中心区域的分部积分与
式~\eqref{eq:artificial-interfaces} 的法向约定。由此定义连续非线性算子函数
\begin{align}
  F_\beta(k):H^1_\beta(\mathcal{C}_0)
  &\longrightarrow H^1_\beta(\mathcal{C}_0)^*,
  \notag\\
  \langle F_\beta(k)u,v\rangle
  &:=\int_{\mathcal{C}_0}\nabla u\mathbin{\cdot}\nabla\overline v
  -k^2\int_{\mathcal{C}_0}\rho u\overline v
  -\sum_{\varsigma\in\{-,+\}}
   \left\langle
   \Lambda_\varsigma(k)\gamma_\varsigma u,
   \gamma_\varsigma v
   \right\rangle_{\Gamma_\varsigma}.
  \label{eq:continuous-center-operator-function}
\end{align}
它关于 $k$ 非线性，因为 $\Lambda_\pm(k)$ 本身依赖于 $k$。

\begin{proposition}[无界导模问题与中心 DtN 问题的等价性]
  \label{prop:continuous-dtn-equivalence}
  固定满足假设~\ref{ass:projected-gap} 的 $\beta$，并令 $k>0$ 满足
  式~\eqref{eq:dtn-admissible-parameters}，且两个半波导 Dirichlet 问题均唯一可解。
  则 $k^2$ 是原无界 $\beta$-表述~\eqref{eq:gap-restricted-beta-formulation} 的导模特征值，
  当且仅当
  \begin{align}
    \ker F_\beta(k)&\ne\{0\}.
    \label{eq:dtn-kernel-criterion}
  \end{align}
  此外，全局导模到 $\mathcal{C}_0$ 的限制给出两个解空间之间的线性双射。
\end{proposition}

\begin{proof}
  设 $u$ 是原无界问题的非零导模。它在 $B_\pm$ 上的限制满足
  问题~\eqref{eq:half-guide-dirichlet-problem}，其 Dirichlet 数据为
  $\varphi_\pm=\gamma_\pm(u|_{\mathcal{C}_0})$。由半波导问题的唯一性，这两个限制正是定义
  $\Lambda_\pm(k)$ 的衰减解，故沿同一法向 $\nu_\pm$ 的法向导数与中心限制相等，并满足
  式~\eqref{eq:center-dtn-problem} 的透明边界条件。分部积分于是给出
  $F_\beta(k)(u|_{\mathcal{C}_0})=0$。若中心限制为零，则其界面迹为零，半波导唯一性又迫使
  两侧限制均为零，这与 $u\ne0$ 矛盾；故 $\ker F_\beta(k)\ne\{0\}$。

  反之，取 $u_0\in\ker F_\beta(k)\setminus\{0\}$。以
  $\varphi_\pm=\gamma_\pm u_0$ 分别求解两个唯一的衰减半波导 Dirichlet 问题，并得到
  $u_-$、$u_+$。弱式~\eqref{eq:continuous-center-operator-function} 先由内部试验函数给出中心区域内的
  分片 Helmholtz 方程和传输条件，再由界面迹的任意性给出
  $\partial_{\nu_\pm}u_0=\Lambda_\pm(k)\gamma_\pm u_0$。因此三部分在
  $\Gamma_\pm$ 上具有相同的 Dirichlet 迹和按同一法向计算的法向导数。将
  $u_-$、$u_0$、$u_+$ 依次拼接，便得到 $B$ 上满足全部内部传输条件和 $y$ 方向
  $\beta$-准周期条件的弱解。式~\eqref{eq:half-guide-decay} 保证拼接后的场属于
  $H^1_\beta(B)$，所以它是原问题的导模；其中心部分非零，故全局场非零。两种构造互为逆，
  从而还得到所述线性双射。
\end{proof}

命题~\ref{prop:continuous-dtn-equivalence} 是针对当前分片常数折射率与传输条件写出的连续
版本；其限制与延拓结构对应于 Fliss 的有界域 DtN 等价定理\cite[定理~4.5]{Fliss2013}。
下一节将构造连续 DtN 算子的数值近似，并据此形成离散非线性特征值问题。
